% =====================================================================
% v2 EXPERIMENTAL PROTOCOL (replaces Sec. IV "Experimental Setup" and the
% protocol paragraphs of Sec. III once the locked runs finish).
% Every value below is taken from ecg_lock/{data,corruption,train,export}.py.
% Result numbers enter ONLY through \input{results_macros.tex}.
% =====================================================================

\subsection{Dataset, Labels, and Splits}
\label{subsec:data}
We use PTB-XL v1.0.3 \cite{wagner2020ptb} at 100~Hz ($T=1000$ samples per 10-s lead). Diagnostic SCP statements are mapped to the five diagnostic superclasses (NORM, MI, STTC, CD, HYP) through \texttt{scp\_statements.csv}; records without any superclass label are excluded, following the PTB-XL benchmark convention \cite{strodthoff2020deep}. The primary analysis uses the recommended stratified folds: folds 1--8 for training, fold~9 for validation, and fold~10 for testing. These folds are patient-disjoint by construction, which our pipeline additionally asserts. As a sensitivity analysis, the primary contrasts are repeated on a patient-level \texttt{GroupShuffleSplit} (80/10/10, seed~7). Split record counts, patient counts, and label prevalences are reported in Table~\ref{tab:split}. Each lead is standardized with the mean and standard deviation computed over all training records.

\subsection{Synthetic Corruption Protocol}
\label{subsec:corruption}
During training, each record is corrupted on the fly as follows. (i)~Additive Gaussian noise is applied to each lead independently with probability $0.45$ and standard deviation $\sigma\sim\mathcal{U}[0,0.25]$ in standardized units. (ii)~With probability $0.5$, $k\sim\mathcal{U}\{1,\dots,6\}$ leads are set to zero (dropout is applied after noise, so dropped leads are exactly zero). (iii)~When the reconstruction loss is active and no lead has been corrupted, one random lead receives noise with $\sigma\sim\mathcal{U}[0.05,0.25]$. (iv)~With probability $0.3$, baseline wander is added to all leads: a sinusoid with amplitude $\mathcal{U}[0.1,0.5]$, frequency $\mathcal{U}[0.05,0.5]$~Hz, and an independent phase per lead. The per-lead target is $r_\ell=1-m_\ell$, where $m_\ell=1$ iff lead $\ell$ received dropout or noise. Baseline wander is deliberately excluded from $m$.

Two protocol variants test whether the per-lead head's clean-input behavior is a training-distribution artifact. The first leaves 20\% of training records fully clean (\texttt{clean-frac} $=0.2$). The second additionally sets a noise floor of $\sigma_{\min}=0.02$ and uses a lead-preserving head (below).

\subsection{Models and Baselines}
The shared encoder, diagnosis heads, pooled per-lead head, and decoder follow Table~\ref{tab:arch}. A \emph{lead-preserving} per-lead head applies three grouped 1D convolutions (12 groups, 16 channels per lead, stride~2) to the input, so that each lead keeps its own feature stream, and maps per-lead mean and standard deviation features to one logit per lead. Diagnosis-only baselines are the same CNN with only the diagnosis head, a Wang-style ResNet1D \cite{wang2017time}, and an InceptionTime-style network \cite{ismail2020inceptiontime}. Every diagnosis-only model is trained twice: on clean inputs, and with the same on-the-fly corruption used for the multitask model. This separates the effect of the auxiliary tasks from the effect of augmentation.

\subsection{Training and Checkpoint Selection}
\label{subsec:training}
All models use AdamW (learning rate $2\times10^{-3}$; $1.5\times10^{-3}$ for the Transformer head; weight decay $10^{-2}$), batch size~128, gradient-norm clipping at~1.0, cosine learning-rate decay, a maximum of 30 epochs, and early stopping with a patience of 8 epochs. The loss weights are $\lambda_{\text{diag}}=1$, $\lambda_{\text{lead}}=0.5$, and $\lambda_{\text{rec}}=1$; inactive tasks have weight~0. The checkpoint is the epoch with the highest macro AUROC on the \emph{clean} validation set (ties: earliest epoch). This rule was fixed before any test evaluation. Each configuration is trained with three seeds (0, 1, 2) using deterministic cuDNN kernels.

\subsection{Evaluation Conditions}
\label{subsec:eval_conditions}
All predictions are exported once per checkpoint as immutable, hashed bundles. Validation bundles are always written before test bundles. The conditions are:
\begin{itemize}
    \item \textbf{Clean} test inputs.
    \item \textbf{Training-distribution corruption}: one fixed-seed realization of the training corruption process, applied to every validation or test record. Per-lead synthetic-corruption-label metrics, reconstruction, repair, and abstention are evaluated here.
    \item \textbf{Stress grids}: exactly $k\in\{0,1,2,3,4,6,8,10,12\}$ dropped leads; Gaussian noise $\sigma\in\{0.05,0.10,0.15,0.20,0.25,0.50\}$ on all leads; baseline wander with amplitude $\{0.25,0.5,1.0\}$; and a joint grid $k\in\{0,2,4,6,8\}\times\sigma\in\{0,0.10,0.25\}$. Each level uses five realizations whose seeds depend only on the condition, so every model sees \emph{identical} corrupted inputs and model contrasts are paired. Levels with $k>6$ or $\sigma=0.50$ lie outside the training distribution.
    \item \textbf{Recorded artifacts}: electrode-motion, muscle, and baseline-wander noise from the MIT-BIH Noise Stress Test Database, resampled to 100~Hz, added to 1, 3, or 6 random leads at 12, 6, and 0~dB SNR.
\end{itemize}

\subsection{Per-Lead Label Baselines, Repair, and Abstention}
The learned synthetic-intact score is compared with transparent rules computed from the corrupted input alone: within-lead standard deviation (flat leads), the first-difference-to-signal energy ratio (additive broadband noise), and their label-free combination (maximum percentile rank). For repair, reconstructed channels replace either the known synthetic mask (\emph{oracle-mask upper bound}) or the leads whose synthetic-intact score falls below $\tau_{\text{mask}}$ (\emph{predicted mask}). $\tau_{\text{mask}}$ maximizes Youden's $J$ on validation and is applied unchanged to the test set. Abstention is evaluated by the area under the risk--coverage curve (AURC), with selective risk defined as the label-wise error rate at $t_{\text{cls}}=0.5$. Accepting records in order of the mean synthetic-intact score is compared with ordering by MC-dropout uncertainty (20 passes; mean over classes of the predictive standard deviation), by the maximum-probability margin, and at random.

\subsection{Statistical Analysis}
Clean-test metrics are reported as mean $\pm$ SD over the three seeds. Contrasts between two configurations are reported as the seed-averaged difference, with a two-sided 95\% patient-level paired percentile bootstrap interval (1{,}000 resamples; patients resampled jointly for both models and all seeds). ECE is the micro label-instance ECE with 15 equal-width bins. All numbers in this paper are generated by one analysis script from the hashed bundles; the bundle hashes, configurations, and code commit are provided in the supplementary material.
